\section{Application: sticky priors for hidden Markov models}\label{sec:sticky}

This section is an application for readers from statistics. Nothing earlier
depends on it, and a reader who is interested only in the general theory and
the planar walk can skip to Section~\ref{sec:verify}.

In a hidden Markov model the hidden states $X_1,\dots,X_N$ form a Markov chain
on $d$ states. Before any data are seen, the model already says how the $N$
time steps are likely to be shared among the states. This is the law $P_N$ of
the count vector $n=(n_1,\dots,n_d)$, where $n_i$ is the number of time steps in
state $i$ (the composition of Section~\ref{sec:model}). A sticky prior makes
the chain leave its current state rarely, which is the regime of this paper.
We ask which states the most likely count vector uses.

Here is what we use from the earlier sections. Suppose the chain leaves its
state with probability $h=\tau\log N/N$ at each step and then picks the next
state from a kernel $\Pi$. For a set $F$ of states (a face) let $\rho_F$ be
the Perron root of $\Pi$ restricted to $F$. The chain stays in $F$ for all $N$
steps with probability about $N^{-\tau(1-\rho_F)}$, and about $N^{|F|-1}$ count
vectors use exactly the states of $F$. Theorem~\ref{thm:first} says that for
large $N$ the most likely count vector uses a set $F$ that minimizes
$\Phi_\tau(F)=\tau(1-\rho_F)+|F|-1$, and that its fractions $n/N$ are then
close to the law $\alpha^*_F$, the entrywise product of the left and right
Perron vectors of $\Pi$ on $F$. Besides this we use the description of a tie
between several minimizers (Theorem~\ref{thm:window}) and the exact identity
for a symmetric pair of states (Proposition~\ref{prop:pair}), once each.

The answer is as follows. As the stickiness decreases, that is, as $\tau$
grows, the most likely count vector moves from one state to larger sets of
states. Among symmetric kernels only the uniform one jumps from one state
straight to all of them.

Fox et al.~\cite{fox2011} add a self-transition bias $\kappa\ge0$ to the HDP
hidden Markov model (HDP stands for hierarchical Dirichlet process), so that
the hidden chain stays in each state longer. In their weak-limit approximation
with base weights $\bar\beta$ and concentration $\omega$, row $j$ of the
transition matrix is Dirichlet with parameters
$\omega\bar\beta_1,\dots,\omega\bar\beta_j+\kappa,\dots,\omega\bar\beta_d$. By
neutrality, row $j$ off the diagonal, divided by $1-\pi_{jj}$, is Dirichlet
with parameters $(\omega\bar\beta_l)_{l\ne j}$ and independent of $\pi_{jj}$. So
$\kappa$ only sets how sticky the chain is, and the kernel of the switches does
not depend on it. In this section we fix every exit probability $1-\pi_{jj}$ at
the same value $h=\tau L/N$, the scale of Kagey's crossing, and study the
kernel. The prior itself mixes over all values of the exit probabilities, and
we do not treat this mixture.

So let $\Pi=(\pi_{ij})$ be a stochastic matrix on $[d]$ with zero diagonal,
$\tau>0$ and $h=\tau L/N<1$. The hidden chain has transition matrix
$(1-h)I+h\Pi=I+hQ$ with $Q=\Pi-I$, so
$q_i=1$ and $T=\tau L$ (Example~\ref{ex:models}(iii)). We assume that
$\pi_{ij}>0$ for all $i\ne j$, which we call (A+), and that the start $\mu$ has
full support. Under the prior (A+) holds almost surely. Let
$\rho_F=\varrho(\Pi_F)$ for the principal submatrix $\Pi_F$ on $F$. Then the
exit rate of $F$ is $\lambda_F=1-\rho_F$, and
$Q_F$ and $\Pi_F$ have the same Perron vectors. A vertex (a single state) has
$\rho_F=0$, and
$\rho_{[d]}=1$. If $2\le|F|\le d-1$, then $\Pi_F$ is irreducible by (A+), and
$0<\rho_F<1$ since $\rho_F$ is at least the least row sum of $\Pi_F$ and the
Perron root of the irreducible $\Pi$ exceeds that of each proper principal
submatrix~\cite{bermanplemmons1994}. So
$\Phi_\tau(F)=\tau(1-\rho_F)+|F|-1$, and we put
\[
 \tau_-=\min_{|F|\ge2}\frac{|F|-1}{\rho_F},\qquad
 \tau_+=\max_{F\ne[d]}\frac{d-|F|}{1-\rho_F}.
\]
We call a global mode of $P_N$, that is, a most likely value of the count
vector before any data are seen, a \emph{prior mode}.

\begin{corollary}[support and location of the prior mode]\label{cor:map}
Assume \textup{(A+)}, let $\mu$ have full support, and fix $\tau>0$. Let
$\mathcal F_\tau$ be the set of faces that minimize $\Phi_\tau$. There is
$N_0(\tau)$ such that for $N\ge N_0(\tau)$ every prior mode $n^*$
satisfies $\supp n^*\in\mathcal F_\tau$ and $n^*/N=\alpha^*_F+O(1/L)$ with
$F=\supp n^*$. If $\mathcal F_\tau=\{F^*\}$, then
$n^*/N\to\alpha^*_{F^*}=l\circ r$, where $l$ and $r$ are the left and right
Perron vectors of $\Pi_{F^*}$ with $l\cdot r=1$. No reversibility is needed.
\end{corollary}

In words, for large $N$ a prior mode uses exactly the states of a set that
minimizes $\Phi_\tau$, and it shares the $N$ time steps among them in the
proportions $\alpha^*_F$.

\begin{proof}
By (A+) every face is irreducible and, as $\mu$ has full support,
admissible. So Theorem~\ref{thm:first}(iii) and~(iv) give the claims, with
$\alpha^*_F=l\circ r$.
\end{proof}

The threshold $N_0(\tau)$ is not uniform in $\tau$. At $\tau=\tau_-$ the
vertices tie with the faces that attain $\tau_-$, and
Theorem~\ref{thm:window}(iii) with $\tau^*=\tau_-$ and a fixed large $t=t_0$
shows that for large $N$ the prior mode is not a single state at
$\tau=\tau_-(1-\frac{\log L-2t_0}{2L})<\tau_-$. So $N_0(\tau)\to\infty$ as
$\tau\uparrow\tau_-$.

The next proposition says for which $\tau$ the prior mode uses one state, all
states, or a set in between.

\begin{proposition}[thresholds]\label{prop:thresholds}
Assume \textup{(A+)} and $d\ge3$.
\begin{enumerate}[(a)]
\item The minimizers of $\Phi_\tau$ are exactly the vertices if and only if
$\tau<\tau_-$, and no vertex minimizes $\Phi_\tau$ if $\tau>\tau_-$. The face
$[d]$ is the unique minimizer if and only if $\tau>\tau_+$, and $[d]$ does not
minimize $\Phi_\tau$ if $\tau<\tau_+$.
\item $1<\tau_-\le d-1\le\tau_+$.
\item $\tau_-=d-1$ if and only if $\tau_+=d-1$, and both hold if and only if
\begin{equation}\label{eq:directjump}
 \rho_F\le\frac{|F|-1}{d-1}\qquad\text{for every face } F.
\end{equation}
If \eqref{eq:directjump} fails, then $\tau_-<d-1<\tau_+$, and for
$\tau\in(\tau_-,\tau_+)$ every minimizer $F$ of $\Phi_\tau$ has
$2\le|F|\le d-1$.
\end{enumerate}
Consequently, let $\mu$ have full support, fix $\tau$ and let $N\ge N_0(\tau)$
as in Corollary~\ref{cor:map}. If $\tau<\tau_-$, the prior
modes are the $Ne_i$ with $\mu_i=\max_l\mu_l$. If $\tau>\tau_+$, every
prior mode has full support. If $\tau_-<\tau<\tau_+$, every prior
mode has a support $F$ with $2\le|F|\le d-1$, and $F=F^*$ if $F^*$ is
the unique minimizer of $\Phi_\tau$.
\end{proposition}

\begin{proof}
We compare every face with a vertex and with $[d]$. The function
$\Phi_\tau(F)=\tau(1-\rho_F)+|F|-1$ is affine in $\tau$, so each comparison
changes sign at one value of $\tau$, and $\tau_-$ and $\tau_+$ are the extreme
ones.

(a) A vertex has $\Phi_\tau=\tau$. For $|F|\ge2$ we have $\rho_F>0$ and
\[
 \Phi_\tau(F)-\tau=|F|-1-\tau\rho_F=\rho_F\Bigl(\frac{|F|-1}{\rho_F}-\tau\Bigr).
\]
So $F$ is worse than a vertex exactly when $\tau<(|F|-1)/\rho_F$. This holds
for every $F$ with $|F|\ge2$ exactly when $\tau<\tau_-$, and if $\tau>\tau_-$
some $F$ is better than every vertex. In the same way $\Phi_\tau([d])=d-1$,
and for $F\ne[d]$ we have $1-\rho_F>0$ and
\[
 \Phi_\tau(F)-(d-1)=\tau(1-\rho_F)-(d-|F|)
 =(1-\rho_F)\Bigl(\tau-\frac{d-|F|}{1-\rho_F}\Bigr).
\]
So $F$ is worse than $[d]$ exactly when $\tau>(d-|F|)/(1-\rho_F)$, which gives
the two claims on $[d]$.

(b) In the minimum that defines $\tau_-$, a face with $2\le|F|\le d-1$ has
$\rho_F<1$ and so $(|F|-1)/\rho_F>1$, and $[d]$ gives the value $d-1\ge2$. So
$\tau_->1$ and $\tau_-\le d-1$. In the maximum that defines $\tau_+$, a vertex
gives the value $d-1$, so $\tau_+\ge d-1$.

(c) By (b), $\tau_-=d-1$ exactly when $(|F|-1)/\rho_F\ge d-1$ for all
$|F|\ge2$, and $\tau_+=d-1$ exactly when $(d-|F|)/(1-\rho_F)\le d-1$ for all
$F\ne[d]$. After clearing the denominators both conditions say
$(d-1)\rho_F\le|F|-1$. This inequality always holds for the vertices and for
$[d]$, so each condition is \eqref{eq:directjump}. If
\eqref{eq:directjump} fails, then neither holds, so $\tau_-<d-1<\tau_+$ by
(b), and (a) excludes the vertices and $[d]$ on $(\tau_-,\tau_+)$.

For the consequence, Corollary~\ref{cor:map} says that the support of a prior
mode minimizes $\Phi_\tau$, and (a) and (c) say which faces can do so. If
$\tau<\tau_-$ the prior modes are vertex compositions, and
$P_N(Ne_i)=\mu_i(1-h)^{N-1}$ is largest for the largest $\mu_i$. Parts (a) to
(c) are checked in Lean.
\end{proof}

We say that $\Pi$ has a \emph{direct jump} when \eqref{eq:directjump} holds.
Then at first order the prior mode goes from one state straight to all $d$ states at
$\tau=d-1$. For the uniform kernel $\pi_{ij}=1/(d-1)$ the matrix $\Pi_F$ has
constant row sums, so $\rho_F=(|F|-1)/(d-1)$. Hence
\eqref{eq:directjump} says that no set of states is more closed than under the
uniform kernel. That kernel has $Q=(J-dI)/(d-1)$, and all its faces tie at
$\tau=d-1$. This is the degeneracy of~\cite{paperB} (Corollary~\ref{cor:paperB}).
The next corollary shows that a direct jump is special.

\begin{corollary}[only the uniform kernel]\label{cor:uniform-kernel}
Let $d\ge3$ and assume \textup{(A+)}. \textup{(a)} If $\Pi$ is symmetric,
then $\Pi$ has a direct jump if and only if $\pi_{ij}=1/(d-1)$ for all
$i\ne j$. \textup{(b)} If $d=3$, then $\Pi$ has a direct jump if and only if
$\pi_{ij}\pi_{ji}\le\frac14$ for all three pairs. In particular every cyclic
kernel $\pi_{i,i+1}=a$, $\pi_{i,i-1}=1-a$ (indices mod $3$) has a direct
jump.
\end{corollary}

\begin{proof}
We test \eqref{eq:directjump} on the pairs. The matrix $\Pi_{\{i,j\}}$ has
zero diagonal and off-diagonal entries $\pi_{ij}$ and $\pi_{ji}$, so
$\rho_{\{i,j\}}=\sqrt{\pi_{ij}\pi_{ji}}$.

(a) Here $\rho_{\{i,j\}}=\pi_{ij}$, so for a pair \eqref{eq:directjump} reads
$\pi_{ij}\le1/(d-1)$. The $d$ rows sum to $1$ and have $d(d-1)$ off-diagonal
entries, so the average of the $\pi_{ij}$ is $1/(d-1)$. If none of them
exceeds the average, they are all equal to it. So a direct jump forces the
uniform kernel. Conversely the uniform kernel has $\rho_F=(|F|-1)/(d-1)$ for
every face, as we saw above, so it satisfies \eqref{eq:directjump} with
equality.

(b) For $d=3$ the faces are the vertices, the three pairs and $[3]$. For the
vertices and $[3]$, \eqref{eq:directjump} reads $0\le0$ and $1\le1$. For a
pair it reads $\sqrt{\pi_{ij}\pi_{ji}}\le\frac12$. The cyclic kernel has
$\pi_{ij}\pi_{ji}=a(1-a)\le\frac14$ for every pair. Both parts are checked in
Lean.
\end{proof}

So for a symmetric kernel that is not uniform there is a range of $\tau$ in
which the prior mode uses more than one state but not all of them
(Proposition~\ref{prop:thresholds}(c)). Without symmetry a direct jump is
possible, and for three states we can say how likely it is under the prior.

\begin{proposition}[the direct jump under the prior]\label{prop:directjump}
Let $d=3$ and let the base weights be uniform. In the weak-limit sticky
HDP-HMM the rows of $\Pi$ are then independent, and the first off-diagonal
entry of each row has the law $\mathrm{Beta}(\omega/3,\omega/3)$. Write
$\Prob_\omega$ for this law of $\Pi$ and $D$ for the event that $\Pi$ has a
direct jump. Then \textup{(a)} $\Prob_\omega(D)\to\frac14$ as $\omega\to0$,
and \textup{(b)} $\Prob_\omega(D)=\frac{3\sqrt3}{8\pi\omega}(1+o(1))$ as
$\omega\to\infty$.
\end{proposition}

\begin{proof}
Let $B_1,B_2,B_3$ be independent with law $\mathrm{Beta}(\omega/3,\omega/3)$,
and put $\pi_{12}=B_1$, $\pi_{21}=B_2$ and $\pi_{31}=B_3$, so that the other
off-diagonal entries are $1-B_j$. By Corollary~\ref{cor:uniform-kernel}(b),
$D$ is the event that $B_1B_2$, $(1-B_1)B_3$ and $(1-B_2)(1-B_3)$ are all at
most $\frac14$.

(a) As $\omega\to0$ each $B_j$ tends in law to a fair coin on $\{0,1\}$. So
the law of $(B_1,B_2,B_3)$ tends to the uniform law on the eight points of
$\{0,1\}^3$. At these points the three products take only the values $0$ and
$1$, so none of the points lies on the boundary of $D$ (where some product
equals $\frac14$). Hence $\Prob_\omega(D)$ tends to the limit probability of
$D$, which is the fraction of the eight points at which all three products are
$0$. Only $(1,0,1)$ and $(0,1,0)$ have this property (checked in Lean), so the
limit is $\frac28=\frac14$. Part (b) is proved in Appendix~\ref{app:hdp}.
\end{proof}

In words, a direct jump has probability close to $\frac14$ when the
concentration is small and of order $1/\omega$ when it is large. For (b), for
large $\omega$ the set $D$ is a thin tube around the cyclic kernels, with
cross-section of area about $\frac98t^4$ at the logit coordinate $t$ along the
cyclic kernels (Appendix~\ref{app:hdp}), which gives the order $1/\omega$. So
the uniform tie is not typical. At $\omega=1$, Monte Carlo with $10^6$ draws
gives $\Prob_\omega(D)=0.115$. So about $0.88$ of random kernels have a range
of $\tau$ in which, at first order, the prior mode sits on a pair
(Proposition~\ref{prop:thresholds}(c)).

The comparison of a single state with a symmetric pair is exact at every $N$.

\begin{corollary}[an exact link to Kagey's crossing]\label{cor:sticky-pair}
Let $i\ne j$ satisfy $\pi_{ij}=\pi_{ji}=\pi_0>0$ and $\mu_i=\mu_j>0$. We do not
assume \textup{(A+)}. Then for every $N\ge2$ the balanced compositions are the
most likely ones with support $\{i,j\}$, and they are more likely than $Ne_i$
if and only if $\tau>\tau^{\rm pair}(N):=Nu_N/((\pi_0+u_N)L)$. If also
$\mu_i=\mu_j=\max_l\mu_l$, then for every $\tau$ with
$\tau^{\rm pair}(N)<\tau<N/L$ no prior mode is a single state. Moreover,
\[
 \tau^{\rm pair}(N)=\frac1{\pi_0}\Bigl[1-\frac{\frac12\log L-c_0}{L}
 +O\Bigl(\frac{\log L}{L^2}\Bigr)\Bigr],
\]
which tends to the first-order value $1/\rho_{\{i,j\}}=1/\pi_0$.
\end{corollary}

In words, a single state against a symmetric pair is Kagey's board, so the
value of $\tau$ at which the pair overtakes the state is known at every $N$ and
not only in the limit.

\begin{proof}
We apply Proposition~\ref{prop:pair}. Every exit rate is $1$, so its
hypotheses hold with $a=\pi_0$ and $q=1$, and its odds are $u=\pi_0h/(1-h)$. It
says that the balanced compositions are the most likely ones on $\{i,j\}$ and
that they beat $Ne_i$ exactly when $u>u_N$. Solving for $h$, this is
$h>u_N/(\pi_0+u_N)$, and with $h=\tau L/N$ it is $\tau>\tau^{\rm pair}(N)$.

For the prior modes, $P_N(Ne_l)=\mu_l(1-h)^{N-1}\le P_N(Ne_i)$ for every $l$.
So when the pair beats $Ne_i$ it beats every single state. The condition
$\tau<N/L$ is $h<1$, and the range of $\tau$ is not empty as
$u_N/(\pi_0+u_N)<1$.

Finally $Nu_N=z_N/p_N$ and $u_N=O(L/N)$, so
$\tau^{\rm pair}(N)=\frac{z_N}{\pi_0L}(1+O(L/N))$, and the second-order
expansion of $z_N$ in~\cite[Corollary~3.7(b)]{paperA} gives the claim.
\end{proof}

So far we looked at the single most likely count vector. The last result
gives the total probability that the chain uses only the states of a set $F$.

\begin{proposition}[face masses]\label{prop:mass}
Assume \textup{(A+)} and $0<h<1$. For every face $F$ with $|F|\ge2$ and
$\Lambda_F=1-h\lambda_F$,
\[
 \frac{\min r_F}{\max r_F}\mu(F)\Lambda_F^{N-1}
 \le\Prob(\supp n\subseteq F)\le
 \frac{\max r_F}{\min r_F}\mu(F)\Lambda_F^{N-1},
\]
and for fixed $\tau$, $\Lambda_F^{N-1}=N^{-\tau\lambda_F}e^{O(L^2/N)}$. Also
$\Prob(\text{one state})=(1-h)^{N-1}\sim N^{-\tau}$ and
$\Prob(\supp n=[d])\to1$ for fixed $\tau$.
\end{proposition}

In words, the chain stays inside $F$ with probability
$\mu(F)N^{-\tau\lambda_F}$ up to a bounded factor, and the probability that it
uses every state tends to $1$.

\begin{proof}
We write the probability with the transition matrix restricted to $F$ and
compare the vector $\ones$ with the Perron vector. Let
$U_F=(1-h)I+h\Pi_F\ge0$. The chain stays in $F$ for all $N$ steps with
probability $\Prob(\supp n\subseteq F)=\mu_F^{\top}U_F^{N-1}\ones$, and
$U_Fr_F=\Lambda_Fr_F$. As $r_F/\max r_F\le\ones\le r_F/\min r_F$ entrywise and
$U_F^{N-1}$ has nonnegative entries, we may replace $\ones$ by these two
multiples of $r_F$. This gives the bounds. For fixed $\tau$,
$(N-1)\log(1-h\lambda_F)=-\tau\lambda_FL+O(L^2/N)$. A vertex composition is a
single path that never switches, which gives $\Prob(\text{one state})$. Every
$F\ne[d]$ has $\lambda_F>0$, so each of the finitely many proper faces has
probability tending to $0$, and $\Prob(\supp n=[d])\to1$. The finite bounds are
checked in Lean.
\end{proof}

\begin{remark}[prior mode versus mass]\label{rem:mapmass}
Let $d\ge3$. For fixed $\tau<\tau_+$ and large $N$ every prior mode has a
proper support, while $\Prob(\supp n=[d])\to1$. For $\tau<\tau_-$ the prior mode uses
one state, an event of probability about $N^{-\tau}$. A face of size $k$ holds
about $N^{k-1}$ compositions, so it can carry most of the mass while each of
them is less likely than a vertex. Table~\ref{tab:sticky} shows this. At
$\tau=0.8$ the prior mode is a single state, yet $\Prob(\text{one state})=0.0122$. At
$\tau=8$ first
order puts the prior mode on $\{1,2\}$, but at $N=240$ it has full support. So the
prior mode can be a poor summary of a sticky prior, and face selection tells us
on which states it sits.
\end{remark}

\begin{table}[t]
\centering\small
\begin{tabular}{@{}rlccc@{}}
\toprule
$\tau$ & prior mode & one state & $\{1,2\}$ & $[3]$\\
\midrule
0.5 & one state & 0.0643 & 0.465 & 0.379\\
0.8 & one state & 0.0122 & 0.423 & 0.531\\
3 & $(120,120,0)$ & $4.3\times10^{-8}$ & 0.129 & 0.871\\
8 & $(111,111,18)$ & $1.2\times10^{-21}$ & 0.0081 & 0.992\\
\bottomrule
\end{tabular}
\caption{The kernel $\Pi$ with rows $(0,0.9,0.1)$, $(0.9,0,0.1)$,
$(0.5,0.5,0)$, the uniform start and $N=240$, where $\tau_-=10/9$ and
$\tau_+=10$. The prior mode (from exact dynamic programming, with a tie of
three compositions at $\tau=0.5$ and $0.8$) and the prior probabilities of the supports.}
\label{tab:sticky}
\end{table}
